%% 第五章 反馈驱动记忆修正

\chapter{反馈驱动记忆修正}
\chaptermark{反馈驱动记忆修正}

\section{引言}

第三章的图增强检索解决了"找得更准"的问题，第四章的效用引导遗忘解决了"存得更省"的问题。然而，长期部署中的记忆系统还面临第三个维护维度的挑战：\textbf{错误的持久化与传播}。在层级记忆树的架构中，每个节点的摘要文本由底层原始数据按需计算得到，本身并非可直接赋值修改的普通字段。视觉场景图构建中的误识别（如将 “cabinet（橱柜）” 误标为 “fridge（冰箱）”）、语音识别中的转写错误、或 LLM 摘要中的概括偏差一旦发生，就会被"冻结"在记忆树中，并通过层级引用关系持续污染所有涉及该节点的后续问答。更严重的是，由于高层摘要递归引用了低层的内容，一个低层节点的错误会被逐层放大，形成从点污染到面污染的传播效应。

认知科学中的\textbf{记忆再巩固}（memory reconsolidation）理论提供了另一种启发：已巩固的长时记忆在被重新激活（回忆或使用）后会短暂地回到不稳定状态，此时它可以被修改、更新或强化，然后再次稳定下来。这一机制使人类记忆在保持整体稳定性的同时，具备了动态修正的能力——记忆不是"写一次就冻结"，而是"每次使用都是一次潜在更新的机会"。当前层级记忆树的一个核心局限恰恰在于缺乏这一机制：记忆条目一旦生成便固化，即使后续发现了错误也无法在不破坏原始数据的前提下进行修正。 GRAF-Mem 的修正模块正针对上述局限而提出——它赋予机器人记忆一种工程化的"再巩固"能力，使系统能够在收到反馈后定位、修正并传播修复。

本章解决第一章所提出的第三个研究问题（RQ3）：如何使记忆系统在收到反馈后能够定位并修正错误，抑制错误在后续问答中的传播。核心方法是\textbf{反馈驱动记忆修正机制}（Feedback-Driven Memory Correction），包含两个层面的创新：（1）\textbf{四阶段修正管线}——反馈锚定的逆向错误定位（事实差异+时间邻近+锚点覆盖+结构源头先验）→ LLM 辅助最小化修正 → 传播检测（时间邻居相似度）→ 自动传播；（2）\textbf{非侵入式视图重定向机制}（Non-intrusive View Redirection / Memory Masking）——在不侵入底层数据的前提下，通过视图重定向覆盖摘要的展示与检索行为，确保修正可逆、可追溯且不破坏原始感知记录。在工程实现上，本文通过 Python 运行时动态属性完成该机制的挂载。

本章 5.2 节详细阐述四阶段管线的算法设计，包括 5.2.1 节反馈锚定的逆向定位（含垂直链扩展机制）、 5.2.2 节 LLM 辅助最小化修正（含先验证后修正的两步式提示词（Prompt）与语义事实判别设计）、 5.2.3 节传播检测（含水平时间邻居与垂直树结构两个传播维度）与 5.2.4 节自动传播修正（含 LLM 类比推理与差异词对回退的双路径策略）。 5.3 节介绍非侵入式视图重定向机制的设计原理及其在语义搜索和树展示两个调用端的集成方式。 5.4 节介绍模拟反馈评测协议——同一 episode 内共享 \texttt{history} 的有状态问答闭环。 5.5 节在受控错误注入范式下，系统验证五个维度：错误定位的垂直全链覆盖能力（5.5.1）、 LLM 修正的语义精确性与判别准确性（5.5.2）、水平与垂直双维度传播检测（5.5.3）、端到端闭环的答案改善（5.5.4）以及修正持久性与工程鲁棒性（5.5.5）。 5.6 节进行综合讨论。

\section{四阶段修正管线}

\subsection{阶段一：错误定位}

阶段一旨在从大规模记忆树中提取与被纠正事实高度相关的候选节点，为后续语义验证与最小化修正提供输入。输入为用户反馈四元组 $(q, a_{\mathrm{wrong}}, a_{\mathrm{correct}}, t_q)$，输出为局部候选集 $C=\{n_1,\ldots,n_m\}$。其中 $C$ 应同时满足三个条件：（1）在时间上位于提问时刻之前；（2）在语义上围绕同一被纠正事实；（3）在结构上覆盖可能的感知源头节点及其上层传播摘要节点。

采用这一设计，是因为纯语义相似度检索存在三类结构性不足。第一，语义近邻并不等于因果相关节点；在长期记忆树中，大量节点可能共享相似对象词或动作词，但它们并不属于同一事件上下文。第二，纯语义检索无法利用时间因果约束，而任何会影响当前回答的错误节点都必须满足 $t_n \le t_q$。第三，用户反馈已经显式给出了被纠正事实的内容，因此阶段一以该事实差异为约束，逆向构造候选范围，而无需复现智能体此前的完整读取路径。

为此，首先从 $(q, a_{\mathrm{wrong}}, a_{\mathrm{correct}})$ 中抽取一个纠错事实四元组

\begin{equation}
z = (e, r, v_{\mathrm{wrong}}, v_{\mathrm{correct}})
\label{eq:5-1}
\end{equation}

其中 $e$ 表示核心实体或事件目标，$r$ 表示被纠正的关系或属性槽位，$v_{\mathrm{wrong}}$ 与 $v_{\mathrm{correct}}$ 分别表示错误值和正确值。例如，对问题 Where did you retrieve the bread from? 及答案对 fridge/cabinet ，可抽取得到 $z=(\mathrm{bread}, \mathrm{retrieved\_from}, \mathrm{fridge}, \mathrm{cabinet})$。该四元组是阶段一与阶段二之间的第一层接口：阶段一据此构造候选集，阶段二据此验证候选节点是否真的承载了被纠正事实。

在得到 $z$ 后，构造三类探针文本：

\begin{equation}
p_{\mathrm{topic}}   = \mathrm{text}(e, r, q)
\label{eq:5-2}
\end{equation}
\begin{equation}
p_{\mathrm{wrong}}   = \mathrm{text}(e, r, v_{\mathrm{wrong}})
\label{eq:5-3}
\end{equation}
\begin{equation}
p_{\mathrm{correct}} = \mathrm{text}(e, r, v_{\mathrm{correct}})
\label{eq:5-4}
\end{equation}

其中 $p_{\mathrm{topic}}$ 用于刻画事件主题，$p_{\mathrm{wrong}}$ 与 $p_{\mathrm{correct}}$ 则分别刻画被纠正事实的错误版本与正确版本。阶段一首先在满足 $t_n \le t_q$ 的节点中，以 $p_{\mathrm{topic}}$ 执行一次粗检索，得到与当前事件主题最相关的一组种子节点；随后沿这些种子节点的父子链路向上、向下各扩展一跳，并补入时间邻近节点，构成局部候选集 $C$。这一扩展步骤的必要性在于：真实错误既可能存在于 L2 原始事件节点，也可能存在于 L3/L4+ 摘要节点，甚至同时以"感知源头 + 摘要传播"的形式分布在一条层级链路上。因此，阶段一输出的不应是单个全局 top-1 节点，而应是一个覆盖相关局部结构的候选集。

在候选集 $C$ 上，进一步定义节点嫌疑度函数

\begin{equation}
\mathrm{suspicion}(n) = \alpha \cdot S_{\mathrm{fact}}(n) + \beta \cdot S_{\mathrm{temp}}(n) + \gamma \cdot S_{\mathrm{anchor}}(n) + \delta \cdot S_{\mathrm{struct}}(n)
\label{eq:5-5}
\end{equation}

其中 $\alpha=0.35$、$\beta=0.30$、$\gamma=0.20$、$\delta=0.15$。这四项信号共同用于对候选集内部节点建立验证优先级。

第一项为\textbf{事实冲突分} $S_{\mathrm{fact}}(n)$。它衡量节点内容更接近错误事实还是正确事实：

\begin{equation}
S_{\mathrm{fact}}(n) = \sigma\big( \mathrm{sim\_vec}(n, p_{\mathrm{wrong}}) - \mathrm{sim\_vec}(n, p_{\mathrm{correct}}) \big)
\label{eq:5-6}
\end{equation}

其中 $\mathrm{sim\_vec}(\cdot,\cdot)$ 为文本嵌入余弦相似度，$\sigma(x)=1/(1+e^{-10x})$ 为 sigmoid 缩放函数。与普通语义检索不同，$S_{\mathrm{fact}}$ 的比较对象不是"节点与问题是否相关"，而是"节点更支持错误版本还是更支持正确版本"。因此，它直接服务于纠错任务本身。

第二项为\textbf{时间邻近度} $S_{\mathrm{temp}}(n)$：

\begin{equation}
S_{\mathrm{temp}}(n) = \exp\big( -|t_q - t_n| / \tau \big)
\label{eq:5-7}
\end{equation}

其中 $\tau$ 为时间衰减常数，默认取 86400 秒。该项先用 $t_n \le t_q$ 施加硬因果约束，再在满足约束的节点中对时间上更接近相关事件的节点给予更高权重。它的作用不是取代语义信息，而是把语义检索限制在因果上合理的局部窗口内。

第三项为\textbf{锚点覆盖分} $S_{\mathrm{anchor}}(n)$。定义锚点集合 $A_z=\{e,r,v_{\mathrm{wrong}},v_{\mathrm{correct}}\}$，并衡量候选节点对这些锚点的覆盖程度：

\begin{equation}
\begin{aligned}
S_{\mathrm{anchor}}(n) = \big( w_e \cdot m(n,e) + w_r \cdot m(n,r) + w_v \cdot \max(m(n,v_{\mathrm{wrong}}), m(n,v_{\mathrm{correct}})) \big) \\
/ (w_e + w_r + w_v)
\end{aligned}
\label{eq:5-8}
\end{equation}

其中 $m(n,\cdot)$ 表示节点文本对指定锚点的匹配强度，默认 $w_e=0.4$、$w_r=0.4$、$w_v=0.2$。该项的作用是保证候选节点落在同一事件框架中，例如区分"另一个无关的 fridge 事件"与"同一个 bread retrieval 事件中的错误位置属性"。

第四项为\textbf{结构源头先验} $S_{\mathrm{struct}}(n)$：

\begin{equation}
S_{\mathrm{struct}}(n) = \lambda \cdot S_{\mathrm{level}}(n) + (1 - \lambda) \cdot \max_{m \in N_{\mathrm{tree}}(n)} S_{\mathrm{fact}}(m)
\label{eq:5-9}
\end{equation}

其中 $S_{\mathrm{level}}(n)=1.0/0.7/0.4$ 分别对应 L2/L3/L4+ 节点，$N_{\mathrm{tree}}(n)$ 表示 $n$ 的父子邻域，$\lambda=0.6$。前一项体现源头优先的层级偏好，后一项体现局部结构一致性：若某条父子链上的多个节点共同支持同一错误事实，则该链路上的节点整体更值得进入阶段二验证。

综合式（5-6）至（5-9），阶段一输出的并不是"最终错误节点"，而是一个按优先级排序的局部候选集 $C$。该候选集将作为阶段二的直接输入。由此，阶段一与阶段二之间的数据流转关系可以表述为：

\begin{equation}
(q, a_{\mathrm{wrong}}, a_{\mathrm{correct}}, t_q) \rightarrow C = \{n_1, \ldots, n_m\}
\end{equation}

其中 $C$ 中的每个节点都附带其嫌疑度及局部上下文信息。阶段一的成败标准也不在于是否把真实污染源排在全树第 1 位，而在于是否能把真实错误节点或其紧邻传播节点稳定纳入一个足够小、足够相关的候选范围，使阶段二能够在该局部范围内完成进一步验证与修正。

\subsection{阶段二：LLM 辅助最小化修正}

阶段二接收阶段一输出的局部候选集 $C$，首先对候选节点进行语义事实判别（\textbf{候选验证}），再对判定为错误事实载体的节点执行\textbf{最小化修正}。因此，阶段二的输出是一个已确认修正的节点集合

\begin{equation}
R = \{r_1, \ldots, r_s\}, \quad R \subseteq C
\end{equation}

其中 $R$ 将作为阶段三传播检测的输入。这样，前两阶段之间的数据流转关系可以表述为：

\begin{equation}
C = \{n_1, \ldots, n_m\} \rightarrow R = \{r_1, \ldots, r_s\}
\end{equation}

这一接口设计的关键在于：阶段一负责高召回地提出候选，阶段二负责在候选内部完成高精度甄别。换言之，阶段一不必独立承担最终裁决责任，而阶段二也不应被理解为"对 top-k 节点直接重写"；它首先要判断某个候选节点是否真的承载了被纠正事实，只有判断为"是"时，才对该节点生成修正后的摘要。

为实现这一过程，修正 LLM 接收如下输入：
（1）候选节点的当前有效摘要；
（2）阶段一输出的纠错事实四元组 $z=(e,r,v_{\mathrm{wrong}},v_{\mathrm{correct}})$；
（3）原始问题 $q$；
（4）错误答案 $a_{\mathrm{wrong}}$；
（5）正确答案 $a_{\mathrm{correct}}$；
（6）该节点在候选局部结构中的上下文信息，如父节点摘要、直接子节点摘要或时间邻近节点摘要。
这些信息共同定义了一个"局部判别问题"：该节点是否处于被纠正事实所在的结构位置上，以及其摘要中是否包含需要替换的错误片段。

因此，阶段二的提示词应遵循两步式目标。第一步，要求 LLM 根据候选摘要及其上下文判断：该节点是否表达了 $e$-$r$-$v_{\mathrm{wrong}}$/$v_{\mathrm{correct}}$ 这一被纠正关系。第二步，仅当答案为肯定时，再按照\textbf{最小化修正}原则生成修正后的摘要文本，即只修改与错误事实直接相关的片段，保留其他信息不变，不添加解释性语句。这样设计的原因有三点。第一，候选集 $C$ 中可能同时包含真实错误节点、传播节点以及仅语义相近但并未出错的邻近节点，因而阶段二必须先做验证再做改写。第二，最小化修正可以降低 LLM 大面积重写摘要所带来的二次偏差。第三，尽量保持文本其余部分不变，有助于阶段三从"修正前/修正后差异"中提取稳定的错误指纹。

当 LLM 调用不可用或失败时，系统回退到\textbf{简单文本替换}策略。该回退策略仍然遵循与阶段二相同的接口：只有当候选节点摘要中确实出现了与 $a_{\mathrm{wrong}}$ 相对应的错误短语时，才执行替换；否则跳过该节点。具体做法是，优先尝试在原始摘要中完整匹配 $a_{\mathrm{wrong}}$ 的关键短语并替换为 $a_{\mathrm{correct}}$ 中的对应部分；若无可匹配的完整片段，则提取两者之间的差异词对并进行局部词级替换。例如，$a_{\mathrm{wrong}} = \text{``I retrieved it from the fridge''}$ 与 $a_{\mathrm{correct}} = \text{``I retrieved it from the cabinet''}$ 可抽取出差异词对 $\mathrm{fridge} \rightarrow \mathrm{cabinet}$，再在候选摘要中进行对应替换。

一旦某个候选节点通过验证并完成修正，系统便将修正后的摘要作为摘要覆盖属性挂载到节点上，同时保留原始摘要副本并清除嵌入缓存。由此，该节点进入已确认修正集合 $R$。只有集合 $R$ 中的节点，才会在阶段三中被视为"修正源节点"，用于继续执行传播检测。这一设计保证了后续传播模块不建立在未经验证的阶段一候选之上，而只建立在已经过阶段二确认的修正结果之上。

在受控注入实验（5.5.2 节）中，阶段二表现出两个重要特征。其一， LLM 能够准确识别摘要中的错误事实并生成对应修正；其二，这种修正不是机械的字符串替换，而是受上下文约束的语义微调。例如，在将 \texttt{"lower fridge"} 修正为 \texttt{"lower cabinet"} 的场景中，LLM 仅替换了错误实体词 \texttt{"fridge"}，而保留了位置修饰语 \texttt{"lower"}，体现了最小化修正仅改变与错误事实直接相关的内容、不重写无关上下文的原则。实验结果表明，阶段二的作用不仅限于文本改写，而是通过先验证后修正的策略将阶段一提出的候选节点收敛为一组语义上已确认、文本上已修正的节点集合 $R$，为后续传播检测提供可靠的修正源。

\subsection{阶段三：传播检测}

错误在记忆中往往不是孤立的。视觉模型在某个时刻的误识别——例如将场景中的 “Toaster（烤面包机）” 误标为 “Microwave（微波炉）”——会同时影响该时刻附近的多个连续感知帧，因为机器人通常在同一场景中持续采集多帧数据。这些连续帧的事件摘要会共享同一个错误的视觉识别结果。因此，从一个已确认修正的节点出发，自动检测其周边可能被同一感知误差污染的其他节点，可以扩大单次修正的覆盖面。

\textbf{水平传播检测}的算法流程如下：（1）从已修正节点获取其修正前的原始错误摘要（即注入错误后、修正前的版本），作为"错误指纹"；（2）将错误指纹编码为嵌入向量；（3）在全局 L2 事件序列中，定位已修正节点的索引位置，对其前后各若干步（默认 2–7 步，取决于预期的污染范围）个时间邻居，计算邻居的有效摘要与错误指纹的余弦相似度；（4）相似度 $\ge$ 阈值（默认 0.7）的邻居标记为疑似传播错误，在其上挂载修正提示元数据（记录传播路径、源节点和置信度）。

传播检测仅跳过同时具有摘要覆盖属性与原始摘要副本的节点——即已完整通过修正管线处理的节点。若一个节点仅有摘要覆盖属性而无原始摘要副本，则表明该节点处于"注入但尚未修正"状态，应被纳入传播检测范围。这一区分避免了将未处理节点的注入错误误判为已修正状态。

\textbf{垂直传播检测}沿树结构而非时间序列进行检测。从已修正节点出发：（1）向上遍历祖先链（父节点→祖父节点→…→根），计算各祖先有效摘要与错误指纹的余弦相似度；（2）向下以 BFS 遍历子孙（子节点→孙节点→…→叶），同样计算相似度；（3）相似度 $\ge$ 阈值（默认 0.65，低于水平传播的 0.7 以补偿跨层级文本风格差异）的祖先/子孙标记为疑似传播错误。向上的覆盖锚点匹配分 $\ge 0.3$ 作为入队条件以避免全树扩散。垂直传播检测与水平传播检测共同构成了从已修正节点出发的双维度传播排查体系。

\subsection{阶段四：自动传播修正}

阶段三检测到的疑似传播错误节点，需通过阶段四进行自动修正，以扩大单次修正的覆盖面。对高置信度（相似度 $\ge 0.8$）的疑似传播错误自动执行类比修正。有 LLM 时，向 LLM 展示"源节点 A 做了从 X 到 Y 的修正，请对相似的邻居节点 B 做相同类型的修正"——LLM 利用类比推理能力将修正映射到语义相似但表述不同的邻居文本。无 LLM 时，从源修正中提取差异词对（如 "Toaster" → "Microwave" → 修正为 "Toaster"），在邻居摘要中执行相同的词级替换。对低置信度（相似度 $< 0.8$）的疑似传播错误，仅挂载修正提示元数据（\texttt{\_correction\_hint}）以待用户确认，不执行自动修正。

\begin{algorithm}[htbp]
\caption{反馈驱动记忆修正管线}
\label{alg:5-1}
\begin{algorithmic}[1]
\Require 问题 $q$，错误答案 $a_{\mathrm{wrong}}$，正确答案 $a_{\mathrm{correct}}$，提问时刻 $t_q$，记忆树 $\mathcal{H}$
\Ensure 修正后的记忆树 $\mathcal{H}'$ 及修正日志 $\mathcal{L}$
\State 从 $(q, a_{\mathrm{wrong}}, a_{\mathrm{correct}})$ 中抽取纠错事实四元组 $z = (e, r, v_{\mathrm{wrong}}, v_{\mathrm{correct}})$ \Comment{式（\ref{eq:5-1}）}
\State 构造主题探针 $p_{\mathrm{topic}}$、错误事实探针 $p_{\mathrm{wrong}}$、正确事实探针 $p_{\mathrm{correct}}$ \Comment{式（\ref{eq:5-2}）–（\ref{eq:5-4}）}
\State 以 $p_{\mathrm{topic}}$ 在 $\{n \in \mathcal{H} \mid t_n \le t_q\}$ 中检索 $M$ 个种子节点，沿父子链路 $\pm 1$ 跳扩展，补入时间邻近节点，执行垂直链扩展，构成候选集 $C$
\For{$\forall n \in C$}
    \State 计算 $\mathrm{suspicion}(n)$ \Comment{式（\ref{eq:5-5}）}
\EndFor
\State 对 $C$ 按嫌疑度降序排列，取 top-$k$ 或按层级 top-$k$ 合并，构成验证集合 $V \subseteq C$（$|V| \le 20$）
\State $R \gets \emptyset$
\For{$\forall n \in V$}
    \State 调用 LLM 判断 $n$ 是否表达 $e$-$r$-$v_{\mathrm{wrong}}$ 错误事实；若 YES，生成修正摘要，执行 \texttt{apply\_summary\_override}，$R \gets R \cup \{n\}$
\EndFor
\For{$\forall n \in R$}
    \State 以 $n$ 修正前摘要为错误指纹，对前后各 $w$ 步时间邻居计算余弦相似度，$\ge 0.7$ 者标记为疑似传播错误 \Comment{水平传播检测}
    \State 沿树结构向上遍历祖先、向下 BFS 遍历子孙，计算余弦相似度，$\ge 0.65$ 者标记为疑似传播错误 \Comment{垂直传播检测}
\EndFor
\For{$\forall$ 疑似传播错误节点}
    \If{相似度 $\ge 0.8$}
        \State 执行类比修正（LLM 类比推理，或回退至差异词对替换）
    \Else
        \State 挂载 \texttt{\_correction\_hint} 元数据以待用户确认
    \EndIf
\EndFor
\State \Return $\mathcal{H}'$ 及修正日志 $\mathcal{L}$（含 $R$、传播检测结果和自动修正记录）
\end{algorithmic}
\end{algorithm}

\section{非侵入式视图重定向机制}

修正操作的最终效果需要通过修改节点的摘要来实现。然而，如前所述，摘要文本是动态计算属性——不可直接赋值。直接修改底层数据（如场景图数据中的物体类别或原始感知数据中的动作指令）会破坏原始感知记录的完整性和一致性。本文提出\textbf{非侵入式视图重定向机制}（Non-intrusive View Redirection，也称 Memory Masking），其设计思路为：不为修正而修改底层数据，而是在数据的展示与检索路径上插入一层可插拔的视图重定向层，使得修正效果对上层调用端可见，而底层感知记录保持完整。

在工程实现上，本文通过 Python 运行时动态属性完成挂载。\texttt{apply\_summary\_override} 函数在目标节点对象上设置三个实例属性：

\begin{itemize}
\item \textbf{摘要覆盖属性（\texttt{\_summary\_override}）}：修正后的摘要文本，是调用端优先读取的修正版本。
\item \textbf{原始摘要副本（\texttt{\_original\_summary}）}：在首次挂载摘要覆盖属性时自动保存的节点原始摘要。该副本支持修正的可逆操作——若需回退，删除摘要覆盖属性即可恢复节点的原始摘要行为。
\item \textbf{修正来源记录（\texttt{\_correction\_source}）}：修正来源的描述字符串（如 \texttt{"injected:object\_swap:cabinet->fridge"} 或 \texttt{"correction"}），用于日志追溯和审计。
\end{itemize}

修正后，节点的嵌入向量缓存被清除，强制语义搜索在下次访问该节点时使用修正后的文本重新计算嵌入向量。

摘要覆盖机制需要在两个主要调用端生效：（1）\textbf{语义搜索端}：语义相似度搜索函数在构建节点的嵌入计算文本时，优先检查摘要覆盖属性的存在——若存在，将覆盖文本追加到索引内容的文本列表中，与子节点的原始摘要一同计算嵌入向量，确保修正后的语义内容被纳入相似度计算。（2）\textbf{交互式树展示端}：\texttt{ExpandableTreeNode.\_\_repr\_\_} 在生成节点的可读摘要时，优先使用摘要覆盖属性（若存在），否则回退到原始的摘要文本，确保智能体在树导航时读取的是修正后的内容。

视图重定向机制的三项设计原则贯穿整个实现：（1）\textbf{非侵入性}：不修改场景图数据、原始感知数据、音频描述等底层字段。摘要覆盖属性等动态属性在 Python 默认的序列化行为下会被序列化，实验验证了它们能完整通过 \texttt{deepcopy} 和 \texttt{pickle} 序列化循环。（2）\textbf{可逆性}：原始摘要副本保留了修正前的版本，任何修正均可通过删除摘要覆盖属性来撤销。（3）\textbf{可追溯性}：修正来源记录保存了每次修正的来源和原因，支持部署后的记忆审计。

\section{模拟反馈评测协议}

标准评测中，每道题从独立的 \texttt{deepcopy(history)} 开始——题目之间完全隔离，前一道题的回答和可能发生的修正不会影响后一道题的记忆状态。这种方式适合评估"从零开始的检索和问答"能力，但无法评估"用户在连续对话中逐步纠正机器人"的真实交互场景。

GRAF-Mem 设计了\textbf{模拟反馈评测协议}（Simulated Feedback Evaluation Protocol），通过模拟反馈评测函数实现。其核心改动在于：同一 episode（即 TEACh 数据中共享同一组 episode ID 的多个问题）内的所有问题\textbf{共享同一份 \texttt{history} 对象}——不进行深拷贝。

具体的评测流程为：（1）对每个新 episode ，将该 episode 的第一个问题对应的 history（深拷贝自预处理缓存）作为共享 \texttt{history}；（2） LLM Agent 使用共享 \texttt{history} 回答当前问题 $Q_i$，生成假设答案 $\mathrm{hypothesis}_i$；（3）答案判别器（Answer Judge）对 $\mathrm{hypothesis}_i$ 与标准答案 $a_i$ 进行两阶段判定——首先进行快速确定性检查（精确字符串匹配、子串包含、 token 序列包含），若均不匹配则调用独立的 Judge LLM 进行语义判定，输出 CORRECT / PARTIAL / WRONG ；（4）若判定为 WRONG ，系统以标准答案 $a_i$ 作为"模拟用户反馈"（相当于用户指出了正确答案），调用完整的四阶段修正管线（5.2 节）修改共享 \texttt{history} 中相关的错误节点；（5）下一道题 $Q_{i+1}$ 在同一份（可能已被 $Q_i$ 的修正管线修改过的）共享 \texttt{history} 上运行。

答案判别器的设计遵循安全优先的保守原则：仅当判定为 WRONG 时触发修正， PARTIAL 判定不触发。答案判别器的提示词明确指示"当在 CORRECT 和 PARTIAL 之间不确定时，倾向 PARTIAL"。这一约束使得系统在答案正确性不确定时不会执行修正，以降低对长期记忆的误改风险。

在系统评测中，共享 \texttt{history} 协议本身（即使修正管线从未被自然触发，详见 5.5 节）导致了与标准评测之间 40\% 的答案差异。这表明有状态反馈闭环对系统行为具有实质性影响：共享的树展开状态、缓存的检索结果和跨题保持的 LLM Agent 上下文，共同构成了记忆状态持续性对后续查询行为的系统性影响。这一现象从侧面验证了反馈闭环的核心前提——记忆的持续性本身就显著改变系统的问答行为，在此之上叠加主动修正，其影响会进一步放大。

为区分"共享状态"与"主动修正"两类不同来源的影响，表 \ref{tab:5-0} 按是否共享 \texttt{history} 与是否触发修正两个维度划分了四种评测配置。组 A 为第五章之前的各章实验所采用的标准独立评测；组 B 仅引入状态持续性、不触发修正，其与组 A 的差异可用于量化持续性效应；组 C 在共享状态基础上叠加保守答案判别器触发的自然反馈闭环；组 D 以确定性规则绕过答案判别器，直接验证修正管线能力。

\begin{table}[htbp]
\centering
\caption{模拟反馈评测的四种配置对比}
\label{tab:5-0}
\begin{tabularx}{\linewidth}{lccX}
\toprule
组别 & 是否共享 \texttt{history} & 是否触发修正 & 目的 \\
\midrule
A（标准评测） & 否 & 否 & 测量原始独立问答性能 \\
B（共享状态） & 是 & 否 & 量化状态持续性对下游查询的影响 \\
C（自然反馈闭环） & 是 & 是，仅 WRONG & 评估保守答案判别器下的自然反馈闭环 \\
D（直接用户反馈） & 是 & 是，确定性触发 & 验证修正管线各阶段的组件能力 \\
\bottomrule
\end{tabularx}
\end{table}

上述 40\% 的答案差异源于组 A 与组 B 之间的对比——即状态持续性效应，而非修正效果。第五章受控实验的主体内容（5.5 节）主要在组 D 配置下展开，以在完全可控的条件下系统验证修正管线。

\section{受控错误注入验证实验}

在自然评测中，修正管线极少被触发——96 次 答案判别器判定中 0 次为 WRONG（85 次 PARTIAL 、 11 次子串/精确匹配）， 20 次判定同样为 0 次 WRONG 。这一现象源于 GRAF-Mem 的两重保护机制：（1）答案判别器的保守判定策略仅在答案明确错误时才触发修正；（2）图增强检索（第三章）的多源证据冗余使孤立节点的语义错误难以穿透智能体综合多源证据后的最终答案。然而，这种自然条件下的零触发使得修正管线各阶段的有效性无法通过常规评测进行定量验证。

因此，本节采用\textbf{受控错误注入范式}：手动向记忆树的特定节点注入已知类型和位置的错误，在完全可控的条件下系统验证修正管线的五个维度：错误定位的垂直全链覆盖能力（5.5.1）、 LLM 修正的语义精确性与判别准确性（5.5.2）、水平与垂直双维度传播检测（5.5.3）、端到端闭环的答案改善链路（5.5.4）以及修正持久性与工程鲁棒性（5.5.5）。

\subsection{错误定位验证}

\textbf{验证思路。}本实验评估反馈锚定错误定位机制在层级记忆树中的覆盖能力。根据 5.2.1 节的管线设计，阶段一输出局部候选集 $C$，阶段二在 $C$ 内执行候选验证与最小化修正，并生成确认修正集 $R$。因此，本实验从两个方面验证定位效果：其一，真实错误节点及其传播邻域是否能够被纳入候选集 $C$；其二，进入候选集的目标节点能否经阶段二验证后形成有效修正集 $R$。

\textbf{实验设置。}在 $|h|=50$ 的预处理记忆树中，在同一垂直链上同时向 L2、L3 和 L4+ 三层节点注入 $\mathrm{cabinet} \rightarrow \mathrm{fridge}$ 错误。具体包括：L2 层原始事件节点（视觉观察中包含 “Cabinet\_2”）、L3 层片段摘要节点、以及 L4+ 高层摘要节点（摘要中描述"从橱柜取出面包"），三处错误均为将 “cabinet/Cabinet” 替换为 “fridge/Fridge”。L2 注入模拟最原始的感知错误，L3 和 L4+ 注入模拟该错误沿层级摘要向上传播后的结果。构造测试 QA = \texttt{"Where did you retrieve the bread from?"}，错误/正确答案分别为 “fridge” 和 “cabinet”，提问时间戳 $t_q$ 设为注入节点时间后约 2 小时。阶段一采用 $M=5$ 的主题种子检索规模，阶段二在候选集内部执行节点验证与最小化修正。

\textbf{垂直链扩展设置。}为验证阶段一对完整垂直链（L2→L3→L4+）的覆盖能力，本实验进一步引入垂直链扩展模块。具体步骤为：对每个已进入候选集的节点，沿树结构向上遍历至根节点、向下以深度优先搜索沿锚点匹配的分支遍历至叶节点（$\mathrm{max\_depth}=8$），将链路上锚点匹配分 $\ge 0.3$ 的节点纳入 $C$。该步骤在式（5-5）的节点嫌疑度排序之前执行。

在相同的注入设置与问题构造下（$M=8$），分别在不启用和启用垂直链扩展两种配置下运行阶段一，结果如表 \ref{tab:5-1} 所示。

\begin{table}[htbp]
\centering
\caption{阶段一候选覆盖结果（含垂直链扩展对比）}
\label{tab:5-1}
\begin{tabular}{lcc}
\toprule
指标 & 不启用垂直链扩展 & 启用垂直链扩展 \\
\midrule
阶段一候选集大小 $|C|$ & 25 & 1221 \\
L2 注入节点是否进入 $C$ & \textbf{否} & \textbf{是} (rank=661) \\
L3 注入节点是否进入 $C$ & 是 (rank=10) & 是 (rank=1101) \\
L4+ 注入节点是否进入 $C$ & 是 (rank=8) & 是 (rank=951) \\
是否进入前 20 验证预算 & L3/L4+ 是 & 否（本实验） \\
\bottomrule
\end{tabular}
\end{table}

\textbf{结果解释。}表 \ref{tab:5-1} 显示，在不启用垂直链扩展的配置下，阶段一的候选集覆盖范围限于 L3/L4+ 层节点（L2 注入节点未进入 $C$），其中 L3/L4+ 节点均处于嫌疑度前 20 名，可在阶段二的自动验证预算内被处理。启用垂直链扩展后，阶段一能够将 L2/L3/L4+ 相关节点全部纳入候选集（L2 $\mathrm{rank}=661$、L3 $\mathrm{rank}=1101$、L4+ $\mathrm{rank}=951$），但三层节点的排名均远超 $\mathrm{max\_verification\_attempts}=20$ 的自动验证预算，无法在完整自动管线中被阶段二逐一验证。

为进一步隔离阶段二的语义修正能力，本文对上述三层目标节点执行定点验证（绕过排名限制）。结果表明，L2、L3 和 L4+ 三类节点均可被阶段二正确验证并修正（均进入 $R$）。由此，配合垂直链扩展机制，阶段一可将候选集覆盖范围从 L3/L4+ 层延伸至包含 L2 感知源节点的完整垂直链；但完整自动管线仍需通过分层 top-k 采样（每个深度层级至少保留 top-k 个候选节点进入 LLM 验证，最终验证集合由全局 top-k 与层级 top-k 合并得到）、链路级重排或锚点阈值调节，使关键节点进入实际验证预算。

垂直链扩展使 $|C|$ 从 25 增至 1221 。虽然 1221 仍显著小于全树 6232 个节点（压缩比约 5:1），但候选集规模的扩大对阶段二的逐一验证提出了效率要求。阶段二设置了 $\mathrm{max\_verification\_attempts}=20$ 的限制：LLM 仅对嫌疑度排名前 20 的候选节点逐一验证，超出此限制则停止。在本实验中，注入节点在启用垂直链扩展后的 $C$ 中排名为 661–1101 ，超出前 20 的范围，因此实验采用了直接对目标节点执行验证的策略，以单独测定管线的修正能力。在实际部署中，可通过分层采样（每个深度层级取 top-k）或提高锚点匹配阈值来控制候选集规模与覆盖范围的平衡。

\subsection{修正质量验证}

\textbf{实验设计。}在 $|h|=50$ 的预处理记忆树中，选取一个 $\mathrm{depth}=4$ 的 L4+高层片段摘要节点节点（LLM 生成的任务级摘要），其原始摘要文本为 \texttt{"I located and retrieved a loaf of bread from the lower cabinet and placed it on the countertop."}。注入物体混淆类型错误——将摘要中的 “cabinet”（橱柜）替换为 “fridge”（冰箱），通过直接设置摘要覆盖属性完成注入（绕过标准的摘要覆盖保存流程，避免自动创建原始摘要副本，从而精确模拟"视觉场景图构建时将橱柜误识别为冰箱"的场景）。构造配套的测试 QA 三元组：问题 = "What did you retrieve the bread from?"，错误答案 = "I retrieved it from the fridge"，正确答案 = "I retrieved it from the cabinet"。调用 LLM 辅助修正流程执行修正。

\textbf{修正质量结果。}在该受控样例中，LLM 成功将摘要中的错误实体 “fridge” 修正为 “cabinet”，并保留了 “lower” 这一位置修饰语，输出 “from the lower cabinet”。这表明该提示词能够在本样例中实现局部、最小化的事实修正，仅改变与错误事实直接相关的实体词，而未重写无关的上下文信息（如位置修饰语 “lower”）。

\textbf{判别准确性结果。}阶段二的提示词设计需处理一类典型歧义场景：在候选集 $C$ 包含大量 L2 感知节点时，这些节点的视觉观察文本中可能自然出现与错误值相同的物体名称（例如场景中真实存在一台冰箱 "Fridge"），但该节点并不表达被纠正的错误主张——它仅将该物体列为环境中的实体，而非将其与实体和关系组合为错误事实。为此，阶段二的提示词要求 LLM 区分两个层次：（1）节点文本中是否\textbf{出现}了错误值词；（2）节点是否\textbf{表达了}将错误值词与被纠正关系结合的特定错误主张。提示词的判别问题表述为："Does this node specifically express the INCORRECT CLAIM that \texttt{<entity>} was \texttt{<relation>} from/in \texttt{<wrong\_value>}? Or does it merely mention \texttt{<wrong\_value>} in an unrelated context?" 实验验证表明，该提示词能够准确识别仅列出环境物体的 L2 节点（返回 NO），并对真正表述错误事实的注入节点（返回 YES），使阶段二的验证建立在语义事实判别而非词级匹配之上。

\subsection{传播检测验证（水平 + 垂直双维度）}

本节分别验证传播检测的两个维度：水平传播（沿时间序列检测连续帧中的同源感知误差）和垂直传播（沿树结构检测跨层级的同源错误传播）。

\textbf{水平传播实验设计。}传播检测的设计假设为连续时间相邻帧受同源感知误差的影响，实验注入场景需与此假设匹配。具体设置为：在同一 Goal 下选取 22 个连续的 L2 事件帧（这些帧共享相同的视觉场景，Toaster 在所有 22 帧中持续出现），对帧 3–7（共 5 帧）注入相同的感知错误——将视觉观察文本中的 “Toaster” 替换为 “Microwave”，模拟视觉模型在连续采集过程中将烤面包机持续误识别为微波炉的场景。以帧 3 为已修正的源节点（手动设置原始摘要副本保存注入的错误版本，并设置摘要覆盖属性保存修正后的版本），运行传播检测流程，搜索范围为前后各 7 跳。

传播检测通过节点属性状态来区分处理阶段：仅跳过同时具有摘要覆盖属性与原始摘要副本的节点（即已完整通过修正管线的节点），而不跳过仅有摘要覆盖属性的节点（即注入但尚未修正的状态）。这一区分保证了未处理的注入节点能够被传播检测覆盖。

\textbf{结果。}传播检测以 100\% 的召回率（4/4）发现了所有 4 个被注入的连续帧（帧 4–7），各帧与源错误指纹的余弦相似度分别为 0.767、0.812、0.962 和 0.756。精确率为 28.6\%（14 个检测结果中 4 个为真阳性），10 个假阳性全部来自同一 Goal 中未被注入的相邻帧——这些帧与注入帧共享几乎完全相同的视觉观察（Toaster 和 Bread\_1\_Sliced\_N 等物体在所有帧中都出现），因此与错误指纹的语义相似度较高（0.71–0.95）。

\begin{table}[htbp]
\centering
\caption{水平传播检测中的真阳性结果}
\label{tab:5-2}
\begin{tabular}{lcc}
\toprule
检测帧（与源帧的距离） & 相似度 & 是否为注入节点 \\
\midrule
帧 4 （+1, Noop） & 0.767 & 是（真阳性） \\
帧 5 （+2, Say） & 0.812 & 是（真阳性） \\
帧 6 （+3, TurnRight） & 0.962 & 是（真阳性） \\
帧 7 （+4, PanLeft） & 0.756 & 是（真阳性） \\
\bottomrule
\end{tabular}

\smallskip
\footnotesize 注：表中仅列出 4 个真阳性检测结果；另有 10 个假阳性来自同一 Goal 下共享视觉场景的未注入帧（详见正文假阳性分析）。
\end{table}

\begin{table}[htbp]
\centering
\caption{传播检测汇总指标}
\label{tab:5-2b}
\begin{tabular}{lc}
\toprule
汇总指标 & 数值 \\
\midrule
召回率（Recall） & \textbf{4/4 = 100.0\%} \\
精确率（Precision） & 4/14 = 28.6\% \\
F1 & 0.444 \\
\bottomrule
\end{tabular}
\end{table}

\textbf{假阳性分析。}10 个假阳性全部来自与源节点共享相同视觉观察的未注入帧。在本实验样例中，假阳性节点原本正确标注了 “Toaster” 且未在其他上下文中包含 “Microwave”，因此自动修正等价于空操作，未观察到额外风险。但这一安全性依赖于具体样例——在开放场景中，假阳性节点可能在其他上下文中真实包含被替换的目标词，此时误修正可能引入新错误。因此，高置信自动传播应配合关系槽位匹配、实体一致性检查或人工确认队列加以保护。传播检测在设计上偏向"宁可多报，不可漏报"——漏检一个真实的污染节点意味着该节点的错误将持续影响未来的问答，而检测层的假阳性可通过上述保护措施在修正层加以控制。

\textbf{自动传播修正。}阶段四的自动类比修正采用双路径策略：当 LLM 可用且返回有效输出时，使用 LLM 的类比推理将源修正模式映射至邻居节点；当 LLM 不可用或返回无效输出时，系统自动回退至差异词对提取与替换——从源修正中提取差异词对（如 $\{\mathrm{Microwave}\} \rightarrow \{\mathrm{Toaster}\}$），在邻居摘要中执行词级替换。在本实验中，自动传播修正了 2/4 帧（帧 5 $\mathrm{sim}=0.812$、帧 6 $\mathrm{sim}=0.962$）；帧 4（$\mathrm{sim}=0.767$）和帧 7（$\mathrm{sim}=0.756$）因相似度低于 0.8 的置信度阈值而未自动修正，其修正提示元数据（\texttt{\_correction\_hint}）被挂载以待用户确认。综上，需区分三个不同层次的指标：\textbf{检测召回率}为 4/4 = 100\%（所有注入帧均被检测到），\textbf{检测精确率}为 4/14 = 28.6\%，\textbf{自动修正覆盖率}为 2/4 = 50\%（仅高置信度帧被自动修正）。

\textbf{垂直传播实验设计。}水平传播检测覆盖了同一层级内的时间相邻节点，但错误在记忆树中还存在另一个传播维度——沿树结构的垂直方向（L2 感知错误 → L3 片段摘要 → L4+ 任务摘要）。为此，新增了垂直传播检测算法（见 5.2.3 节），从已修正节点出发沿树结构向上遍历祖先、向下 BFS 遍历子孙，计算各节点有效摘要与源错误指纹的余弦相似度。

在 5.5.1 节的三层级同源注入场景下（L2/L3/L4+ 均注入 “fridge”），以修正后的 L3 节点为源运行垂直传播检测，结果如表 \ref{tab:5-3} 所示。

\begin{table}[htbp]
\centering
\caption{垂直传播检测结果（L3 注入节点为源）}
\label{tab:5-3}
\begin{tabular}{lccc}
\toprule
方向 & 检测数 & 相似度范围 & 节点类型 \\
\midrule
向下（L3→L2 子孙） & 8 & 0.663–0.746 & EventBasedSummary \\
向上（L3→L4+ 祖先） & 0 & — & — \\
\bottomrule
\end{tabular}
\end{table}

垂直传播检测成功发现了 L3 注入节点下 8 个 L2 事件子孙（相似度 0.66–0.75），这些 L2 节点与注入的 L3 共享同一 Goal ，视觉观察文本高度重叠。未检测到向上传播（L3→L4+）的原因是 L3 与 L4+ 属于不同抽象层级，二者的摘要粒度和表述风格差异较大，导致嵌入相似度低于阈值。

与水平传播相比，垂直传播的相似度整体偏低（0.66–0.75 vs 0.76–0.96）。这反映了层级树中跨级节点的文本表征差异： L3/L4+ 的自然语言摘要与 L2 的感知数据格式（\texttt{"Action: \ldots\ Visual observation: \ldots"}）在嵌入空间中距离较远。这一差异提示未来可以通过多模态错误指纹（如在视觉物体标签层面而非自然语言摘要层面进行匹配）来增强跨层级检测的灵敏度。

\subsection{端到端闭环验证}

前三个实验分别在隔离条件下验证了定位、修正和传播检测的组件功能。端到端闭环验证的目标是确认：从智能体产生错误答案 → 用户提供反馈 → 四阶段修正管线执行 → 智能体在修正后的记忆上重新回答同一问题 → 答案改善，这一完整链路能否在受控条件下正确串联。

\textbf{实验设计。}在 $|h|=50$ 的预处理记忆树中，选取 5.5.2 节所用的同一 L4+ 节点（摘要 \texttt{"I located and retrieved a loaf of bread from the lower cabinet..."}），注入 $\mathrm{cabinet} \rightarrow \mathrm{fridge}$ 错误。构造三组测试问题：两组靶心问题直接命中注入事实（"Where did you retrieve the bread from?"、"What container did you take the bread out of?"），一组对照问题测试无关事实（"Did you use a toaster?"）。对每个问题，以注入后的 \texttt{history} 为上下文调用 LLM Agent 生成回答。

本实验绕过答案判别器的保守触发策略，采用确定性规则触发修正：当智能体回答中包含 “fridge”（不区分大小写）时，直接构造反馈四元组 $(q, a_{\mathrm{wrong}}, a_{\mathrm{correct}}, t_q)$ 并调用修正管线。这一步不改变修正管线的逻辑，只是模拟了"用户明确指出你记错了"的交互场景（5.6 节讨论的触发场景 a）。修正管线执行后，再次用相同问题查询智能体，对比修正前后答案。

\textbf{定位与修正结果。}在不启用垂直链扩展的配置下 $|C|=19$，启用后 $|C|=147$。注入的 L4+ 目标节点在两种配置下均进入 $C$（不启用时 $\mathrm{rank}=9$，启用后 $\mathrm{rank}=118$——启用垂直链扩展后候选集中纳入了更多高嫌疑度的 L2 节点，目标排名相应后移）。阶段二按照 5.5.2 节所述的提示词设计，对该节点执行 LLM 先验证后修正：LLM 判别该节点表达了被纠正事实（VERDICT: YES），并生成了修正摘要 \texttt{"I located and retrieved a loaf of bread from the lower cabinet and placed it on the countertop."}——\texttt{"fridge"} 被替换为 \texttt{"cabinet"}，其余内容不变。

\textbf{答案改善结果。}修正前后的 智能体答案对比如表 \ref{tab:5-4} 所示。

\begin{table}[htbp]
\centering
\small
\caption{用户直接反馈闭环中修正前后答案对比}
\label{tab:5-4}
\begin{tabular}{p{3.2cm}p{4.0cm}p{4.0cm}}
\toprule
问题 & 修正前答案 & 修正后答案 \\
\midrule
Where did you retrieve the bread from?
  & \texttt{"From the lower \textbf{fridge}."}（错误）
  & \texttt{"From the lower \textbf{cabinet}."}（正确） \\
What container did you take the bread out of?
  & \texttt{"The lower \textbf{fridge}."}（错误）
  & \texttt{"I took the bread out of the lower \textbf{cabinet}."}（正确） \\
Did you use a toaster? (对照)
  & \texttt{"Yes, I put a slice of bread in the toaster."}（正确）
  & \texttt{"Yes, I put a slice of bread in the toaster."}（正确，无变化） \\
\bottomrule
\end{tabular}
\end{table}

修正前，两个靶心问题的答案均包含错误值 “fridge”；修正后，两个靶心问题的答案均更正为 “cabinet”，且对照问题的答案保持不变，表明修正未引入副作用。端到端实验在受控条件下验证了从用户反馈输入到 智能体答案改善的完整修正链路：反馈输入 → 阶段一定位 → 阶段二 LLM 验证并修正 → 智能体重新回答 → 答案改善。

\textbf{传播检测。}本次实验为单节点注入（仅 L4+），在同一 episode 的其余节点未注入同源错误。在该设置下，水平传播检测和垂直传播检测均未返回候选传播节点，说明传播检测模块在无同源污染条件下未产生额外误报。

\textbf{与自然评测零触发的关系。}5.5 节开头提到，在自然评测（不注入错误）中，答案判别器的保守设计导致修正管线 0 次触发（96 次判定中 0 次 WRONG）。自然评测中的低触发率表明，修正模块在当前系统中更接近一种兜底维护机制：图增强检索（第三章）的多源证据冗余降低了孤立错误影响最终答案的概率，修正管线则用于在用户明确反馈（本实验所模拟的场景）或系统性感知污染（5.5.3 节模拟的场景）出现时进行定点维护。两者是互补关系，而非互斥。

\subsection{修正持久性验证}

修正管线通过视图重定向机制（5.3 节）实现修正效果，该机制声明了三项设计原则：非侵入性、可逆性、可追溯性。本节对这些工程属性进行正式的实验验证。

\textbf{实验设计。}在 $|h|=50$ 的记忆树中选取一个 L4+ 节点，调用 \texttt{apply\_summary\_override} 应用修正（$\mathrm{cabinet} \rightarrow \mathrm{sideboard}$），然后依次测试六个维度：（1）非侵入性——节点的 \texttt{nl\_summary} 属性（从底层原始数据动态计算的摘要）是否保持原始值；（2）可逆性——删除摘要覆盖属性后是否恢复为原始摘要副本；（3）可追溯性——修正来源记录是否完整记载修正轨迹；（4）\texttt{deepcopy} 存活——三个修正相关属性是否通过 \texttt{copy.deepcopy} 循环完整保留；（5）\texttt{pickle} 存活——上述属性是否通过 \texttt{pickle.dumps/loads} 序列化循环完整保留；（6）索引可检索性——修正后的摘要文本是否能被语义搜索命中。

\textbf{结果。}六项测试全部通过，如表 \ref{tab:5-5} 所示。

\begin{table}[htbp]
\centering
\small
\caption{修正持久性测试结果}
\label{tab:5-5}
\begin{tabular}{p{0.20\linewidth} p{0.15\linewidth} p{0.55\linewidth}}
\toprule
测试维度 & 结果 & 关键证据 \\
\midrule
非侵入性 & 通过 & \texttt{nl\_summary} 保持原始值（含 \texttt{"cabinet"}），\texttt{\_summary\_override} 为 \texttt{"sideboard"} \\
可逆性 & 通过 & 删除 \texttt{\_summary\_override} 后有效摘要恢复为 \texttt{\_original\_summary} \\
可追溯性 & 通过 & \texttt{\_correction\_source} = \texttt{"persistence\_test: cabinet→sideboard"} \\
\texttt{deepcopy} 存活 & 通过 & \texttt{\_summary\_override}、\texttt{\_original\_summary}、\texttt{\_correction\_source} 三者完全保留 \\
\texttt{pickle} 存活 & 通过 & 同上，序列化往返后三者完全保留 \\
索引可检索性 & 通过 & 修正文本自相似度 = 0.894，嵌入后可被语义搜索召回 \\
\bottomrule
\end{tabular}
\end{table}

上述结果验证了 5.3 节提出的三项设计原则。其中，\texttt{pickle} 存活测试确认了修正状态能够完整通过预处理缓存文件的序列化管道，这是 5.5.4 节实验中注入后的 \texttt{history} 能够通过序列化进入评测管道的前提条件。

\section{讨论}

综合 5.5.1–5.5.5 五个受控实验的验证维度，修正模块的实验结果可从以下六个方面加以讨论。

\textbf{（1）反馈锚定逆向定位的层级覆盖范围。}5.5.1 节的实验表明，在 $\pm 1$ 跳父子扩展的配置下，阶段一的候选集覆盖范围限于 L3/L4+ 层节点（L2 注入节点未进入 $C$）。引入垂直链扩展后，阶段一能够将 L2/L3/L4+ 层节点纳入候选集（$|C_{\mathrm{enhanced}}|=1221$，L2 $\mathrm{rank}=661$，L3 $\mathrm{rank}=1101$，L4+ $\mathrm{rank}=951$），但排序靠后。在定点验证策略下，阶段二具备修正这些节点的能力。完整自动管线仍需通过分层 top-k 采样（每层级保留 top-k 候选节点合并验证集合，如 5.5.1 节所述）使关键节点进入实际验证预算。候选集规模从 25 增至 1221 反映了覆盖完整性与搜索聚焦性之间的权衡——1221 仍为全树 6232 个节点的约五分之一，但实际部署中仍需通过分层采样或锚点匹配阈值调节来控制验证开销。

\textbf{（2）阶段二的判别精度。}阶段二的提示词设计区分两个语义层次：节点文本中是否出现了错误值词，以及节点是否表达了将该错误值词与被纠正关系相结合的特定错误主张。这一区分在候选集包含大量自然出现错误值词的节点时尤为关键——例如，L2 感知帧的视觉观察中真实存在的冰箱（“Fridge”）会被正确识别为环境实体而非错误事实载体。5.5.2 节的实验结果表明，该提示词设计使 LLM 能够对仅列出环境物体的节点返回 NO，对真正表述错误主张的注入节点返回 YES，使候选验证建立在语义事实判别而非词级匹配之上。

\textbf{（3）双维度传播检测的覆盖范围与局限。}5.5.3 节的实验覆盖了传播检测的两个维度。水平传播检测在连续帧同源误识别场景下达到 100\% 召回率（4/4），精确率为 28.6\%，自动修正覆盖率为 50\%（2/4）。假阳性集中于与注入帧共享相同视觉场景的未注入帧；在本实验样例中这些假阳性未导致误修正，但开放场景下仍需配合实体一致性检查等保护措施。垂直传播检测从 L3 注入节点出发检测到 8 个 L2 子孙节点（$\mathrm{sim}$ 0.66–0.75），但未检测到向上传播（L3→L4+）。垂直传播的相似度整体低于水平传播（0.66–0.75 vs 0.76–0.96），其原因是跨层级节点间的文本风格差异（L2 的感知数据格式与 L3/L4+ 的自然语言摘要）在嵌入空间中表现为系统性距离，提示多模态错误指纹（如视觉物体标签级匹配）可作为增强跨层级检测灵敏度的方向。

\textbf{（4）端到端修正链路的闭环验证。}5.5.4 节的闭环实验在受控条件下验证了从用户反馈输入到 智能体答案改善的完整链路：修正前 LLM Agent 对靶心问题的答案包含错误值 “fridge”，修正后两个靶心问题的答案均更正为 “cabinet”，且对照问题的答案保持不变。该实验确认了定位、验证修正、传播检测三个阶段在端到端流程中的协同效果：阶段一将目标节点纳入候选集，阶段二完成语义验证与最小化修正，修正后的记忆使 LLM Agent 在重查询时给出正确回答。

\textbf{（5）视图重定向机制的工程可靠性。}5.5.5 节的持久性测试在非侵入性、可逆性、可追溯性三个维度上验证了该机制满足 5.3 节所述的设计规约。具体而言，节点的 \texttt{nl\_summary} 属性保持原始值而不受覆盖影响，删除摘要覆盖属性后有效摘要恢复为原始摘要副本中保存的版本，修正来源记录完整记载修正轨迹；此外，三个覆盖属性在 \texttt{deepcopy} 和 \texttt{pickle} 序列化往返后均完整保留，确保了修正状态能够通过预处理缓存管道进入评测流程。

\textbf{（6）修正触发频率与检索鲁棒性的关系。}自然评测中修正管线的极低触发率（96 次判定 0 次 WRONG）表明，修正模块在当前系统中更接近一种兜底维护机制：图增强检索（第三章）的多源证据冗余降低了孤立错误影响最终答案的概率，修正管线则用于在用户明确反馈（5.5.4 节模拟的场景）或系统性感知污染（5.5.3 节模拟的场景）出现时进行定点维护。检索鲁棒性与修正机制构成互补关系：前者降低修正的触发频率，后者保证触发后的修正质量。修正模块的实际触发场景包括：（a）用户直接提供反馈，指出系统中存在的错误事实，系统以用户提供的 wrong/correct 答案对直接触发修正管线；（b）视觉模型在连续多帧上的系统性感知误识别（如 5.5.3 节模拟的场景），传播检测对受污染节点进行批量发现与修正。

\medskip
\noindent\textbf{本章工作的局限。}综上，本章实验仍存在以下局限。第一，错误定位在启用垂直链扩展后虽然提升了候选覆盖率，但候选集规模显著扩大（$|C|$ 从 25 增至 1221），关键节点未必进入默认验证预算（$\mathrm{max\_verification\_attempts}=20$），后续仍需引入分层 top-k 或链路级重排机制。第二，传播检测在水平连续帧场景下具有较高召回率（100\%），但精确率仍较低（28.6\%），开放场景中需要结合实体一致性检查或人工确认机制降低误修正风险。第三，垂直传播检测受跨层级文本风格差异影响，向上传播（L3$\rightarrow$L4+）的检测能力仍有限。第四，当前实验主要基于受控错误注入，真实用户反馈场景下的触发频率、反馈噪声和长期修正效果仍需进一步验证。
